% =====================================================================
%  PETase-ML — Two-stage, leakage-audited protein-stability framework
%  MDPI-FORMATTED VERSION (extended).
%
%  *** IMPORTANT — READ BEFORE SUBMITTING ***
%  This file reports THE SAME STUDY as main.tex (Springer Nature /
%  Scientific Reports formatting). They are two journal formattings of one
%  piece of work, not two papers. Submit ONE of them. Submitting both is
%  dual submission and is research misconduct.
%
%  This file supersedes the earlier "Condition-Aware Machine Learning
%  Framework ..." MDPI draft, which was based on the older
%  19,071-mutation model. Per main.tex, those results are superseded and
%  must not appear.
%
%  PROVENANCE OF EVERY NUMBER IN THIS FILE:
%    - Headline metrics, dataset counts, leakage counts, model tables:
%      main.tex (authoritative) — verified identical.
%    - Average precision, best-F1 thresholds, confusion matrices:
%      paper_figures/perf_summary.json  + the diagnostics figures.
%      NOTE: Stage 2 average precision is stored as exactly 0.4535 — a tie at
%      the third decimal. It is reported here as 0.453 to match the value
%      rendered in precision_recall_curve.pdf, so text and figure agree.
%    - Extended threshold tables (Appendix C):
%      paper_figures/stage1_threshold_sweep_full.csv (1,122 rows)
%      paper_figures/stage2_threshold_sweep_full.csv (669 rows)
%    - IsPETase case study: paper_figures/ispetase_stage1.json
%      (verified absent from training data, S669, and BRENDA benchmark).
%    - Pinned software: multitask/requirements-lock.txt
%  NOTHING in this file is estimated, rounded from memory, or invented.
%
%  SCOPE (editorial, not printed), inherited from main.tex:
%    - temperature-focused study; pH analysis and EpHod excluded (future
%      work only)
%    - general protein-stability framework evaluated for POTENTIAL
%      PET-degrading-enzyme application (not PETase-specific training data)
%    - no wet-lab / industrial / environmental performance claims
%    - the older 19,071-mutation results must not appear
%
%  YELLOW [TO ADD: ...] markers are genuine gaps to resolve before
%  submission. Remove every marker before final submission.
%
%  BUILD: requires the MDPI LaTeX template (provides Definitions/mdpi.cls),
%  available from https://www.mdpi.com/authors/latex . Copy these figures
%  from paper_figures/ into the project root (or adjust \graphicspath):
%    roc_curve.pdf  precision_vs_threshold.pdf  precision_recall_curve.pdf
%    stage1_diagnostics.pdf  stage2_diagnostics.pdf  stage1_tm_scatter.pdf
%    stage2_model_comparison.pdf  internal_vs_external.pdf
% =====================================================================
\documentclass[journal,article,submit,moreauthors]{Definitions/mdpi}

\usepackage{amsmath}
\usepackage{booktabs}
\usepackage{multirow}
\usepackage{siunitx}
\usepackage{xcolor}

\graphicspath{{./}{paper_figures/}}

% ---- highlighted "to add" markers (remove before final submission) ----
\newcommand{\todo}[1]{\colorbox{yellow}{\textbf{[TO ADD: #1]}}}

% ---- shorthand -------------------------------------------------------
\newcommand{\ddg}{\ensuremath{\Delta\Delta G}}
\newcommand{\Tm}{\ensuremath{T_\mathrm{m}}}
\newcommand{\dtm}{\ensuremath{\Delta T_\mathrm{m}}}
\newcommand{\kcal}{\,kcal\,mol$^{-1}$}

\firstpage{1}
\makeatletter
\setcounter{page}{\@firstpage}
\makeatother
\pubvolume{1}
\issuenum{1}
\articlenumber{0}
\pubyear{2026}
\copyrightyear{2026}

%=================================================================
\Title{A Two-Stage, Leakage-Audited Machine Learning Framework for
Temperature- and Stability-Aware Screening of PET-Degrading Enzymes}

% NOTE: author order and affiliations follow main.tex. main.tex itself
% carries an open question on affiliation 1 (see below) — resolve before
% submission.
\Author{Ayush Iyer\textsuperscript{1,}*, James Ponzio\textsuperscript{2}
and Abhinav Iyer\textsuperscript{3}}

\AuthorNames{Ayush Iyer, James Ponzio and Abhinav Iyer}

% NOTE (from main.tex): confirm whether UC Santa Cruz is Ayush Iyer's
% institutional affiliation or a compute-resource acknowledgement only.
% The previous draft listed "Stanford University", which is not supported
% by any record in this project and has been removed.
\address{%
$^{1}$ \quad University of California, Santa Cruz, Santa Cruz, CA, USA\\
$^{2}$ \quad \todo{James Ponzio --- department, organization, city, state, country}\\
$^{3}$ \quad \todo{Abhinav Iyer --- department, organization, city, state, country}}

\corres{Correspondence: iyer.ayush31@gmail.com}

\abstract{Plastic pollution is a defining environmental challenge, and enzymatic
depolymerization of polyethylene terephthalate (PET) offers a biological route to
recycling; however, candidate enzymes must remain stable under demanding conditions, and
computational screening is frequently inflated by data leakage. We built a general
protein-stability framework and evaluated it for potential application to PET-degrading
enzyme engineering. We assembled 1{,}001{,}888 mutation and stability records from seven
public sources and routed them by measurement type into a two-stage design: Stage~1
screens whole-protein thermostability (melting temperature, \Tm) and Stage~2 optimizes
mutation-level folding stability (\ddg). Training records were audited for both exact and
sequence-homology leakage against the independent benchmarks; auditing removed
approximately 132{,}000 records, of which 131{,}479 were detectable only through sequence
homology and would have survived exact-match filtering entirely. On the independent BRENDA
benchmark, Stage~1 reached a receiver-operating-characteristic (ROC) AUC of 0.732 (95\% CI
0.708--0.755) with an average precision of 0.844 against a positive-class baseline of 0.63.
On the independent S669 benchmark, Stage~2 reached a ROC AUC of 0.669 (95\% CI 0.62--0.72),
an average precision of 0.453 against a baseline of 0.25, and an ensemble Pearson
correlation of 0.390 (95\% CI 0.32--0.46). Both stages showed a pronounced
precision--recall trade-off in which precision rose and recall fell under stricter
thresholds. Stage~1 internal validation (Pearson $\approx$0.80) substantially exceeded
independent performance (0.54), illustrating the size of the optimism that internal
splits alone would have reported. These results describe a moderate, leakage-controlled
screen-then-optimize framework suited to prioritizing candidates for experimental testing,
not a validated industrial solution.}

\keyword{protein stability; machine learning; enzyme engineering; thermostability;
mutation-effect prediction; PET-degrading enzymes; data leakage; protein language models}

\begin{document}

%=================================================================
\section{Introduction}
%=================================================================

\subsection{Plastic Accumulation and the Case for Biological Recycling}

Polyethylene terephthalate (PET) is among the most heavily produced synthetic polymers,
and its accumulation as waste is a major contributor to global plastic pollution: global
plastics production reached roughly 460 million tonnes in 2019 and continues to rise,
while only a small fraction is recycled \citep{oecd2022plastics}.
% CONFIRM: verify the exact OECD figure/year and add any PET-specific value.
Mechanical recycling, the dominant industrial route, progressively shortens polymer chains
and degrades material quality, so recycled PET is typically down-cycled rather than
returned to its original application. Biological recycling, in which enzymes depolymerize
PET into recoverable monomers that can in principle be repolymerized, is an attractive
complement to mechanical and chemical recycling \citep{yoshida2016}.

\subsection{PET Hydrolases and the Role of Thermostability}

Enzymatic PET degradation has advanced through the discovery and engineering of PET
hydrolases such as PETase and related cutinases
\citep{yoshida2016,tournier2020,lu2022}. The discovery of \textit{Ideonella sakaiensis}
201-F6 and its secreted PET hydrolase established that PET is biologically accessible
\citep{yoshida2016}; subsequent protein-engineering campaigns produced variants with
substantially improved operating characteristics \citep{tournier2020,lu2022}.

Thermostability matters in this setting because reaction rates rise near PET's
glass-transition temperature, where the polymer becomes more accessible to enzymatic
attack. It is not, however, the sole determinant of catalytic performance: activity,
substrate binding, reaction conditions, polymer crystallinity, and pH also matter, and a
variant stabilized at the expense of turnover may be worse in practice. We therefore treat
stability as a screening criterion that narrows a candidate set, not as a proxy for overall
enzyme quality.

\subsection{Two Distinct Prediction Problems}

Computational stability prediction addresses two related but non-interchangeable
questions. \emph{Whole-protein screening} asks whether a given enzyme is thermostable,
summarized by its melting temperature (\Tm). \emph{Mutation-level optimization} asks how a
specific substitution changes folding stability, expressed as the change in unfolding free
energy (\ddg). The two differ in their unit of analysis (a protein vs.\ a point mutation),
in the data available to train them, and in the error that matters. Conflating them leads
to misleading claims about what a model can do---for example, reporting a whole-protein
classification accuracy and implying it applies to mutation ranking.

This study keeps the two problems explicitly separate, trains a distinct model for each on
a distinct subset of the data, and evaluates each on its own independent benchmark.

\subsection{Data Leakage in Stability Benchmarking}

A recurring problem in this literature is data leakage. \emph{Exact-match} leakage
(identical sequences or mutations shared between training and test) is straightforward to
detect and is routinely filtered. \emph{Sequence-homology} leakage---where a training
protein is not identical to a test protein but shares a homologous background with
it---is subtler, is not caught by exact-key filtering, and can inflate benchmark
performance substantially because protein stability is strongly conserved within families.

We do not claim that all prior models contain leakage; rather, we control for it explicitly
at both levels, quantify how much each level removes, and report the effect. As shown in
Section~\ref{sec:res-leakage}, homology-level contamination dominated by roughly two orders
of magnitude in this dataset, which has direct implications for how stability benchmarks
should be constructed.

\subsection{Related Work}

Structure-based methods estimate \ddg{} from an explicit three-dimensional model, and
statistical-potential and machine-learning predictors trained on curated thermodynamic
datasets have been developed over two decades; the S2648 dataset and the PoPMuSiC
predictor are representative of this line \citep{dehouck2009s2648}. Such methods require
structures and are sensitive to their quality.

Sequence-based predictors avoid that requirement. Protein language models trained on large
unlabeled sequence corpora---of which the ESM-2 family is a widely used example
\citep{lin2023esm2}---produce residue-level representations that transfer to a range of
downstream property-prediction tasks, including stability, without explicit structural
input. Our framework is sequence-based and uses ESM-2 representations as its sole learned
featurization.

Benchmarking practice has itself become a subject of study. Pancotti~\textit{et al.}
assembled the S669 dataset specifically to provide a test set with limited overlap against
common training sets, and compared a broad set of available tools under matched conditions
\citep{pancotti2022s669}. We adopt S669 as our Stage~2 benchmark for that reason, and we
restrict quantitative comparison to settings where the benchmark, sign convention, split,
and metric genuinely match.

For whole-protein thermostability, large-scale experimental resources have changed what is
trainable: the Meltome Atlas provides proteome-wide thermal stability measurements
\citep{jarzab2020meltome}, the Tsuboyama~\textit{et al.} mega-scale dataset provides
folding stability for hundreds of thousands of designed and natural domains
\citep{tsuboyama2023}, and site-saturation mutagenesis of human domains provides
abundance-based stability proxies at comparable scale \citep{beltran2025domainome}. BRENDA
aggregates curated enzyme functional and stability annotations \citep{chang2021brenda} and
serves here as an independent whole-protein benchmark.

\subsection{Contributions}

Because whole-protein selection and mutation-level optimization are usually treated
separately, this study evaluates a single, integrated two-stage screen-then-optimize
framework---a general protein-stability method assessed for potential application to
PET-degrading enzyme engineering---on leakage-audited, independent BRENDA (temperature) and
S669 (\ddg) benchmarks, and reports its performance without overstating readiness.
Specifically, we contribute:

\begin{enumerate}
\item A harmonized 1{,}001{,}888-record stability table assembled from seven public
      sources and routed by measurement type, so that each stage trains only on the
      measurement it predicts (Section~\ref{sec:sources}, Table~\ref{tab:data}).
\item A two-level leakage audit---exact-key and MMseqs2 sequence-homology---with the
      contribution of each level quantified separately (Section~\ref{sec:leakage},
      Table~\ref{tab:leakage}).
\item Independent evaluation of both stages with bootstrap confidence intervals, reported
      alongside internal validation so that the generalization gap is visible rather than
      concealed (Section~\ref{sec:internal-external}, Figure~\ref{fig:intext}).
\item A full threshold characterization of both stages, showing that high precision is
      obtainable only at low recall, and stating explicitly that benchmark precision is not
      an experimental success rate (Sections~\ref{sec:stage1-thresh} and
      \ref{sec:stage2-class}; Appendix~\ref{app:sweeps}).
\end{enumerate}

%=================================================================
\section{Materials and Methods}
%=================================================================

\subsection{Study Design}

The framework comprises two stages that use different prediction targets and different
data subsets: Stage~1 (whole-protein \Tm) and Stage~2 (mutation-level \ddg). The full
1{,}001{,}888-record table did not train both models; records were routed by measurement
type (Table~\ref{tab:data}, Panel~B). This routing is a deliberate design constraint: a
\Tm{} record describes a protein, a \ddg{} record describes a substitution, and pooling
them would train each model on a target it is not being evaluated against.

\subsection{Source Datasets}
\label{sec:sources}

Records were drawn from seven public sources spanning both measurement types
(Table~\ref{tab:data}, Panel~A): the Human Domainome abundance dataset
\citep{beltran2025domainome}; the Tsuboyama~\textit{et al.} mega-scale folding-stability
dataset and its double-mutant subset \citep{tsuboyama2023}; the Meltome Atlas
\citep{jarzab2020meltome}; ThermoMutDB \citep{xavier2021thermomutdb}; FireProtDB
\citep{stourac2021fireprotdb}; and the S2648 stability dataset \citep{dehouck2009s2648}.
Independent evaluation used the BRENDA enzyme benchmark \citep{chang2021brenda} for
Stage~1 and the S669 benchmark \citep{pancotti2022s669} for Stage~2; protein sequences
were resolved through UniProt.

The Tsuboyama and Domainome data were obtained from their Zenodo deposits (records 7992926
and 13629491, respectively), and the BRENDA temperature-stability labels were parsed from
the BRENDA 2026.1 release. The source datasets were downloaded and assembled between April
and June 2026, and each was used under its respective public license (BRENDA under
CC~BY~4.0).
% CONFIRM: the exact license of each source, and the ProDDG/S2648 provenance
% (S2648 = Dehouck et al.).

\subsection{Harmonization and Sequence Resolution}

Records were harmonized to a common schema, typed by measurement, de-duplicated, and
mapped to canonical UniProt sequences retrieved through the UniProt REST API
(\texttt{rest.uniprot.org}) in June 2026 (training-set sequences on 18~June~2026 and
BRENDA benchmark sequences on 22~June~2026).
% CONFIRM: access dates (from fetch-output file timestamps); add canonical vs. isoform
% handling, failed-mapping count, and dedup policy.

Sign conventions were normalized across sources at load time, because the sources do not
agree: for \ddg{} we adopt the convention that negative values denote stabilization, and
each source was checked against that convention rather than assumed to follow it.
% CONFIRM: state per-source sign handling explicitly (at least one source required
% inversion at load time); reviewers commonly ask for this.

\subsection{Label Definition for Stage 1}
\label{sec:stage1-labels}

BRENDA reports heterogeneous temperature-stability evidence, not a uniform \Tm{} column.
Positive (thermostable) proteins were therefore defined as those with a derived
stable-temperature ceiling $\geq 60\,^{\circ}$C and negative proteins as those
$\leq 45\,^{\circ}$C. Proteins falling in the intervening 45--60\,$^{\circ}$C band were
treated as ambiguous and excluded from classification rather than assigned to either class,
so that the reported classification metrics are not inflated by forcing a decision on
genuinely borderline cases.

For the regression analysis we distinguish two kinds of reported value: a directly measured
melting temperature, and a stability ceiling inferred from inactivation experiments or
qualitative stability evidence. Both are retained, but they are reported separately
(Section~\ref{sec:stage1-reg}) because they are not equally reliable targets.

\subsection{Leakage Auditing}
\label{sec:leakage}

Every candidate training record was audited against the independent benchmark proteins
using (i) an exact-key and mutation-level overlap check and (ii) a sequence-homology search
with MMseqs2 \citep{steinegger2017mmseqs2}. Searches were run with
\texttt{mmseqs easy-search} at maximum sensitivity ($-s\ 7.5$, $-e\ 1000$,
\texttt{-{}-max-seqs 500}), and a training record was flagged as leakage when it matched a
benchmark protein at $\geq$30\% sequence identity over $\geq$50\% query coverage (applied
bidirectionally).

Four independent exact keys were used, so that a record could be caught by any of them:
UniProt accession plus mutation; a hash of the wild-type sequence; PDB identifier plus
chain plus mutation; and PDB identifier plus mutation. Per-key counts are given in
Table~\ref{tab:leakage}.

\subsection{Feature Representation}

Proteins and mutations were represented using ESM-2 protein-language-model embeddings
\citep{lin2023esm2}. Per-residue representations were taken from the final layer of
\texttt{esm2\_t30\_150M\_UR50D} (640 dimensions) and mean-pooled over the sequence.
Sequences longer than the model's 1022-residue context limit were truncated.

For the Stage~1 whole-protein model the feature vector was the mean-pooled embedding
together with the log sequence length. For the Stage~2 mutation model the features combined
the wild-type and mutant mean-pooled embeddings and their difference
($[\mathbf{e}_{wt},\ \mathbf{e}_{mut},\ \mathbf{e}_{mut}-\mathbf{e}_{wt}]$) with an ESM-2
log-likelihood-ratio score for the substitution. These embedding features were reduced by
principal component analysis to 256 dimensions before the boosting models.

Embeddings were precomputed once and cached, so that the language model was evaluated a
single time per sequence rather than once per training epoch.
% CONFIRM: final feature set and PCA dimensionality (GPU job runs build_features.py
% --pca 256); embedding model esm2_t30_150M_UR50D (some earlier scripts use
% esm2_t12_35M); do not carry features from the 19,071-mutation model.

We note a known limitation of this representation: mean-pooling over the sequence
attenuates mutation-site detail, since a single substitution changes only one of several
hundred pooled positions. The explicit difference vector and the log-likelihood-ratio score
partially compensate, but a per-residue or attention-pooled representation is a natural
avenue for improvement (Section~\ref{sec:future}).

\subsection{Models and Hyperparameters}

Each stage used an ensemble that averaged the predictions of its base models, all with a
fixed random seed of 42. Stage~1 (\Tm) combined Extra Trees and Random Forest, XGBoost
\citep{chen2016xgboost}, LightGBM \citep{ke2017lightgbm}, and CatBoost
\citep{prokhorenkova2018catboost}. Stage~2 (\ddg) combined Random Forest, Extra Trees, a
multilayer perceptron, XGBoost, LightGBM, and CatBoost. Random Forest, Extra Trees, and the
multilayer perceptron were implemented in scikit-learn \citep{pedregosa2011sklearn}. The
ensemble prediction was the unweighted mean of the base-model outputs. Complete
hyperparameter settings are given in Appendix~\ref{app:hyper}, Table~\ref{tab:hyper}.

\subsection{Training Infrastructure}

GPU components were trained on the National Research Platform (Nautilus) Kubernetes cluster
using the \texttt{pytorch/pytorch:2.1.0-cuda12.1-cudnn8-runtime} container on the
\texttt{large-gpu} partition, with code, data, and the embedding cache held on a persistent
volume. Sequence-homology auditing used MMseqs2 release~18-8cc5c.

\subsection{Evaluation and Statistical Analysis}

Models were trained on a protein-grouped training/validation split (grouped so that no
protein spanned both partitions; seed 42) and evaluated on the held-out, leakage-audited
BRENDA and S669 benchmarks. Classification was summarized by ROC AUC, average precision,
precision, recall, F1, and accuracy across decision thresholds; regression by Pearson and
Spearman correlation and RMSE.

Confidence intervals were estimated by bootstrapping the test items with replacement
(2{,}000 resamples; percentile 95\% CI; seed 0), reusing the saved per-item predictions
without retraining.

For \ddg, negative values denote stabilization; a variant is predicted stabilizing when its
predicted \ddg{} falls below the decision threshold. Thresholds were swept across the full
range of observed predictions to characterize the precision--recall trade-off rather than
to select an operating point on the benchmark: 1{,}122 distinct thresholds for Stage~1 and
669 for Stage~2 (Appendix~\ref{app:sweeps}).
% CONFIRM: thresholds were fixed before evaluation.

Software versions are listed in Appendix~\ref{app:software}, Table~\ref{tab:software}.
% CONFIRM: these versions were reproduced from the training container (original run
% unpinned); confirm with a pinned re-run.

%=================================================================
\section{Results}
%=================================================================

\subsection{Dataset Composition}

After harmonization, the multi-source staging table comprised 1{,}001{,}888 mutation and
stability records drawn from seven public sources (Table~\ref{tab:data}, Panel~A). Because
the framework treats temperature stability and folding stability as distinct problems,
records were routed by measurement type (Table~\ref{tab:data}, Panel~B): 29{,}654
melting-temperature (\Tm) records were assigned to Stage~1 and 508{,}693 folding
free-energy change (\ddg) records to Stage~2. Abundance-proxy (457{,}943) and \dtm{}
(5{,}598) records were retained in the table but were not used for model training in this
study.

The composition is markedly imbalanced by source: two datasets (Domainome and the
Tsuboyama mega-scale set) supply 81\% of all records, while the three classical curated
thermodynamic databases (ThermoMutDB, FireProtDB, and ProDDG/S2648) together supply about
2\%. The large sources are also the ones furthest from the measurement of interest---the
Domainome records are an abundance-based proxy rather than a direct thermodynamic
measurement, which is why they were excluded from training here.

\begin{table}[H]
\caption{Composition of the assembled dataset and allocation by measurement type. Exact
record counts; both panels sum to 1{,}001{,}888.}
\label{tab:data}
\begin{tabular}{lr}
\toprule
\multicolumn{2}{l}{\textit{Panel A --- Rows by source dataset}}\\
\midrule
\textbf{Source dataset} & \textbf{Rows}\\
\midrule
Domainome (aPCA abundance)              & 457{,}943\\
Tsuboyama \textit{et al.} 2023 (mega-scale)     & 357{,}155\\
Tsuboyama \textit{et al.} 2023 (double mutants) & 138{,}275\\
Meltome Atlas                           & 27{,}884\\
ThermoMutDB                             & 11{,}520\\
FireProtDB                              & 6{,}549\\
ProDDG/S2648                            & 2{,}562\\
\midrule
\textbf{Total}                          & \textbf{1{,}001{,}888}\\
\bottomrule
\end{tabular}

\vspace{8pt}
\begin{tabular}{lrl}
\toprule
\multicolumn{3}{l}{\textit{Panel B --- Rows by measurement type and model allocation}}\\
\midrule
\textbf{Measurement type} & \textbf{Rows} & \textbf{Used to train}\\
\midrule
\ddg{} (mutation)     & 508{,}693 & Stage 2\\
Abundance (proxy)     & 457{,}943 & Not used in this study\\
\Tm{} (whole protein) & 29{,}654  & Stage 1\\
\dtm{}                & 5{,}598   & Not used in this study\\
\midrule
\textbf{Total}        & \textbf{1{,}001{,}888} & ---\\
\bottomrule
\end{tabular}
\end{table}

\subsection{Leakage Audit}
\label{sec:res-leakage}

From 1{,}135{,}644 stability-relevant candidate records, the audit removed 2{,}627 exact
duplicate measurements and 132{,}205 records that matched the S669 test set. Of the
test-set matches, 131{,}536 were flagged by sequence homology and the remainder by exact
keys, with 131{,}479 of the homology matches invisible to every exact key
(Table~\ref{tab:leakage}). This left 1{,}001{,}888 training-eligible records.

The asymmetry is the central observation: the four exact keys together caught 669 records,
while homology search caught 131{,}536---roughly 197 times as many. A pipeline that filtered
only on exact identity would have retained more than 131{,}000 contaminated records, or
about 13\% of the final training table. Because stability is conserved within protein
families, those records would have allowed the model to encounter close relatives of the
test proteins during training, and any resulting benchmark score would overstate
generalization by an unknown but potentially large margin.

\begin{table}[H]
\caption{Leakage audit against the S669 benchmark. Exact keys are applied independently, so
a record may be caught by more than one; the final row reports homology matches that no
exact key caught. Starting pool: 1{,}135{,}644 candidate records; 2{,}627 exact duplicate
measurements were removed separately.}
\label{tab:leakage}
\begin{tabular}{lr}
\toprule
\textbf{Detection route} & \textbf{Records flagged}\\
\midrule
Exact key: UniProt accession $+$ mutation      & 300\\
Exact key: wild-type sequence hash             & 232\\
Exact key: PDB $+$ chain $+$ mutation          & 70\\
Exact key: PDB $+$ mutation                    & 67\\
\midrule
Sequence homology (MMseqs2, $\geq$30\% id, $\geq$50\% cov) & 131{,}536\\
\midrule
\textbf{Total records matching S669}           & \textbf{132{,}205}\\
\quad of which detectable \textit{only} by homology & 131{,}479\\
\midrule
\textbf{Training-eligible records remaining}   & \textbf{1{,}001{,}888}\\
\bottomrule
\end{tabular}
\end{table}

\subsection{Internal Validation Versus Independent Testing}
\label{sec:internal-external}

During development, the Stage~1 temperature model achieved an internal held-out
root-mean-square error (RMSE) of approximately 6.5\,$^{\circ}$C and a validation Pearson
correlation of approximately 0.80. On the independent BRENDA benchmark the same model
obtained a Pearson correlation of 0.54---a drop of roughly 0.26 in correlation
(Figure~\ref{fig:intext}).

We report both figures deliberately. The internal number describes performance on data
drawn from the training distribution and characterizes model fitting only; it is not
evidence of generalization. Had we reported the internal validation alone, as is common,
the framework would have appeared substantially stronger than the independent evaluation
supports. Stage~2 internal validation is not reported here because a comparable
development-split figure was not retained for that stage.

\begin{figure}[H]
\includegraphics[width=0.72\linewidth]{internal_vs_external.pdf}
\caption{Internal (development) validation compared with independent-benchmark performance.
Stage~1 Pearson correlation falls from approximately 0.80 on the internal held-out split to
0.54 on the independent BRENDA benchmark. Stage~2 is shown with its independent S669
correlation (0.39); a comparable internal figure was not retained for that stage.}
\label{fig:intext}
\end{figure}

\subsection{Stage 1: Classification on BRENDA}

Independent evaluation used a curated BRENDA whole-protein temperature-stability benchmark
\citep{chang2021brenda}, labeled as described in Section~\ref{sec:stage1-labels}. During
curation, 412 records were removed from a pre-curation labeled set of 2{,}446 to yield
2{,}034 proteins; a further 471 ambiguous cases in the 45--60\,$^{\circ}$C band were
excluded, leaving 1{,}563 proteins (979 positive, 584 negative) for classification.

On this independent set the classifier reached a ROC AUC of 0.732 (95\% CI 0.708--0.755).
The distribution of predicted \Tm{} separates the two classes only partially
(Figure~\ref{fig:s1diag}, left): both classes concentrate between 45 and 60\,$^{\circ}$C,
and the separation that does exist comes largely from a distinct high-\Tm{} mode near
70--85\,$^{\circ}$C that contains almost exclusively positive proteins. At the best-F1
threshold of 46\,$^{\circ}$C the model recovers 953 of 979 positives but also calls 503 of
584 negatives positive (Figure~\ref{fig:s1diag}, right), giving precision 0.65, recall 0.97,
and F1 0.78.

\begin{figure}[H]
\includegraphics[width=\linewidth]{stage1_diagnostics.pdf}
\caption{Stage~1 classifier diagnostics on the independent BRENDA benchmark
($n=1{,}563$). \textbf{Left:} distribution of predicted \Tm{} by true label; the dashed
line marks the best-F1 threshold (46\,$^{\circ}$C). The classes overlap heavily below
60\,$^{\circ}$C, with a largely positive high-\Tm{} mode near 70--85\,$^{\circ}$C.
\textbf{Right:} confusion matrix at that threshold (best F1 $=0.78$).}
\label{fig:s1diag}
\end{figure}

\subsection{Stage 1: Threshold Behaviour}
\label{sec:stage1-thresh}

Classification performance depended strongly on the decision threshold
(Table~\ref{tab:stage1}, Panel~A; Figure~\ref{fig:precision}a). As the temperature cutoff
was made more stringent, precision increased while recall fell---from 0.65 (recall 0.98) at
the most permissive cutoff to 0.93 (recall 0.39) and 0.98 (recall 0.32) at the strictest
cutoffs. The full sweep across 1{,}122 distinct thresholds is summarized in
Appendix~\ref{app:sweeps}, Table~\ref{tab:sweep1}.

We emphasize that precision at a strict threshold is not an experimental success rate: the
0.93 precision corresponds to only 39\% recall, so the strict-threshold regime identifies a
small, high-confidence set of candidates while missing most true positives. At the
60\,$^{\circ}$C threshold the individual models differ modestly, with Extra Trees the most
precise (0.963) and CatBoost the most sensitive (recall 0.400); the ensemble sits between
them (Table~\ref{tab:stage1}, Panel~B).

\begin{table}[H]
\caption{Stage 1 melting-temperature model: classification on the independent BRENDA
benchmark (1{,}563 enzymes; 979 positive, 584 negative). Positive class $=$ thermostable.
ROC AUC $=0.732$ (95\% CI 0.708--0.755); average precision $=0.844$ against a
positive-class baseline of 0.63. Precision is reported with recall, F1, and accuracy so
that it is not read as overall reliability.}
\label{tab:stage1}
\begin{tabular}{lcccc}
\toprule
\multicolumn{5}{l}{\textit{Panel A --- Decision-threshold analysis}}\\
\midrule
\textbf{Threshold} & \textbf{Precision} & \textbf{Recall} & \textbf{F1} & \textbf{Accuracy}\\
\midrule
46\,$^{\circ}$C & 0.65 & 0.98 & 0.78 & 0.66\\
50\,$^{\circ}$C & 0.71 & 0.77 & 0.74 & 0.66\\
52\,$^{\circ}$C & 0.77 & 0.63 & 0.69 & 0.65\\
60\,$^{\circ}$C & 0.93 & 0.39 & 0.55 & 0.60\\
66\,$^{\circ}$C & 0.98 & 0.32 & 0.49 & 0.58\\
\bottomrule
\end{tabular}

\vspace{8pt}
\begin{tabular}{lccc}
\toprule
\multicolumn{4}{l}{\textit{Panel B --- Per-model performance at the 60\,$^{\circ}$C threshold}}\\
\midrule
\textbf{Model} & \textbf{Precision} & \textbf{Recall} & \textbf{F1}\\
\midrule
Extra Trees   & 0.963 & 0.369 & 0.533\\
Random Forest & 0.948 & 0.355 & 0.517\\
LightGBM      & 0.933 & 0.395 & 0.555\\
CatBoost      & 0.927 & 0.400 & 0.559\\
XGBoost       & 0.917 & 0.394 & 0.551\\
Ensemble      & 0.945 & 0.388 & 0.550\\
\bottomrule
\end{tabular}
\end{table}

\subsection{Stage 1: Regression Against Measured Temperature}
\label{sec:stage1-reg}

Evaluated as a regressor against the reported stable temperature across all 2{,}034 curated
proteins, the model obtained a Pearson correlation of 0.54 (95\% CI 0.50--0.57). Restricted
to the 325 proteins whose reported value is a directly measured melting temperature, rather
than a ceiling inferred from inactivation or qualitative evidence, the correlation rose to
0.60 (Figure~\ref{fig:s1scatter}).

The scatter also exposes a systematic compression: predictions concentrate in a
45--60\,$^{\circ}$C band regardless of the measured value, so proteins with measured
stability above roughly 80\,$^{\circ}$C are consistently under-predicted, and those below
about 35\,$^{\circ}$C are over-predicted. The model is therefore better at ranking proteins
than at recovering absolute temperatures, which is consistent with its intended use as a
screen and is the reason we report classification behaviour as the primary Stage~1 result.

\begin{figure}[H]
\includegraphics[width=0.75\linewidth]{stage1_tm_scatter.pdf}
\caption{Stage~1 predicted versus measured thermostability on the curated BRENDA set
($n=2{,}034$). Grey points ($n=1{,}709$) are proteins whose label is a stability ceiling
inferred from inactivation or qualitative evidence; teal points ($n=325$) are proteins with
a directly measured \Tm. Pearson $r=0.54$ over all points and $r=0.60$ over the measured-\Tm{}
subset. The dashed line is the identity. Predictions are compressed toward
45--60\,$^{\circ}$C relative to the measured range.}
\label{fig:s1scatter}
\end{figure}

\subsection{Stage 2: Regression on S669}

Independent folding-stability performance was assessed on the S669 benchmark of 669 single
mutations (168 stabilizing, 501 non-stabilizing reference cases)
\citep{pancotti2022s669}. Across 2{,}000 bootstrap resamples, the ensemble achieved a
Pearson correlation of 0.390 (95\% CI 0.32--0.46), RMSE 1.509\kcal, and a ROC AUC of 0.669
(95\% CI 0.62--0.72) for stabilizing/non-stabilizing classification (Table~\ref{tab:stage2},
Panels~A--B).

A notable feature of the per-model breakdown is that the ensemble did \emph{not} produce the
best regression metrics: Extra Trees alone reached Pearson 0.436 and RMSE 1.481, ahead of
the ensemble's 0.390 and 1.509 (Figure~\ref{fig:s2models}). Unweighted averaging pulled the
strongest single model toward weaker members---the multilayer perceptron, at Pearson 0.298,
was the weakest---so the ensemble is more conservative but not more accurate on this
benchmark. We report the ensemble as the primary model because it was fixed in advance,
but we note that a weighted or selected ensemble is an obvious improvement and that
choosing Extra Trees post hoc on the basis of test performance would itself be a form of
benchmark overfitting.

\begin{figure}[H]
\includegraphics[width=\linewidth]{stage2_model_comparison.pdf}
\caption{Stage~2 per-model regression performance on the independent S669 benchmark
($n=669$). \textbf{Left:} Pearson correlation (higher is better). \textbf{Right:} RMSE in
\kcal{} (lower is better). The unweighted ensemble (highlighted) ranks second on both
metrics, behind Extra Trees.}
\label{fig:s2models}
\end{figure}

\begin{table}[H]
\caption{Stage 2 mutation-level \ddg{} model: regression and classification on the
independent S669 benchmark (669 mutations; 168 stabilizing, 501 non-stabilizing reference
cases). Ensemble ROC AUC $=0.669$ (95\% CI 0.62--0.72); average precision $=0.453$ against a
positive-class baseline of 0.25; ensemble regression Pearson $=0.390$ (95\% CI 0.32--0.46);
2{,}000 bootstrap resamples.}
\label{tab:stage2}
\begin{tabular}{lccc}
\toprule
\multicolumn{4}{l}{\textit{Panel A --- Regression (predicted vs.\ experimental \ddg)}}\\
\midrule
\textbf{Model} & \textbf{Pearson} & \textbf{Spearman} & \textbf{RMSE (\kcal)}\\
\midrule
Extra Trees           & 0.436 & 0.461 & 1.481\\
XGBoost               & 0.382 & 0.413 & 1.514\\
Random Forest         & 0.369 & 0.396 & 1.523\\
LightGBM              & 0.330 & 0.364 & 1.549\\
CatBoost              & 0.324 & 0.362 & 1.557\\
Multilayer perceptron & 0.298 & 0.338 & 1.612\\
Ensemble              & 0.390 & 0.414 & 1.509\\
\bottomrule
\end{tabular}

\vspace{8pt}
\begin{tabular}{lcccc}
\toprule
\multicolumn{5}{l}{\textit{Panel B --- Ensemble classification across decision thresholds}}\\
\midrule
\textbf{Threshold} & \textbf{Precision} & \textbf{Recall} & \textbf{F1} & \textbf{Accuracy}\\
\midrule
0.00\kcal    & 0.81 & 0.10 & 0.18 & 0.77\\
$+$0.25\kcal & 0.51 & 0.26 & 0.35 & 0.75\\
$+$0.50\kcal & 0.38 & 0.46 & 0.42 & 0.67\\
$+$1.00\kcal & 0.31 & 0.80 & 0.45 & 0.51\\
\bottomrule
\end{tabular}

\vspace{8pt}
\begin{tabular}{lcc}
\toprule
\multicolumn{3}{l}{\textit{Panel C --- Per-model performance at the strict (0.00\kcal) threshold}}\\
\midrule
\textbf{Model} & \textbf{Precision} & \textbf{Recall}\\
\midrule
Random Forest         & 0.900 & 0.107\\
Extra Trees           & 0.786 & 0.065\\
XGBoost               & 0.735 & 0.149\\
LightGBM              & 0.625 & 0.119\\
CatBoost              & 0.429 & 0.107\\
Multilayer perceptron & 0.415 & 0.202\\
Ensemble              & 0.810 & 0.101\\
\bottomrule
\end{tabular}
\end{table}

\subsection{Stage 2: Classification Behaviour}
\label{sec:stage2-class}

As in Stage~1, precision and recall traded off with threshold stringency
(Table~\ref{tab:stage2}, Panels~B--C; Figure~\ref{fig:precision}b). At the strict
0.00\kcal{} threshold the ensemble reached precision 0.81 at only 10\% recall---18 true
stabilizing mutations recovered out of 168, with 4 false positives. The best-F1 operating
point lies far to the permissive side, at $+1.16$\kcal, where recall rises to 0.89 but
precision falls to 0.31 (F1 $=0.46$; Figure~\ref{fig:s2diag}).

The predicted-score distributions show why (Figure~\ref{fig:s2diag}, left): stabilizing and
destabilizing mutations overlap across almost the entire predicted range, with the
stabilizing distribution shifted only modestly toward lower predicted \ddg. The model
therefore operates as a high-precision, low-recall filter at strict cutoffs: the top-ranked
candidates are enriched for true stabilizers, but the majority of stabilizing mutations are
not recovered. We report this low recall explicitly rather than obscuring it, because it
defines the model's intended use as a candidate-prioritization filter.

\begin{figure}[H]
\includegraphics[width=\linewidth]{stage2_diagnostics.pdf}
\caption{Stage~2 classifier diagnostics on the independent S669 benchmark ($n=669$).
\textbf{Left:} distribution of predicted \ddg{} by true label; the dashed line marks the
best-F1 threshold ($+1.16$\kcal). The two classes overlap across most of the predicted
range. \textbf{Right:} confusion matrix at that threshold (best F1 $=0.46$).}
\label{fig:s2diag}
\end{figure}

\subsection{Precision--Recall Behaviour and Average Precision}

Because both benchmarks are class-imbalanced, we also report average precision, which
summarizes the precision--recall curve without reference to a chosen threshold
(Figure~\ref{fig:pr}). Stage~1 attained an average precision of 0.844 against a
positive-class baseline of 0.63, and Stage~2 attained 0.453 against a baseline of 0.25.

Both stages therefore improve on their respective baselines---by 0.21 and 0.20 in absolute
terms---but the shapes differ. The Stage~1 curve holds precision above 0.95 out to roughly
30\% recall before declining smoothly, whereas the Stage~2 curve is unstable at low recall
and falls below 0.6 precision by 20\% recall. The Stage~2 ranking is informative at the very
top of the list and degrades quickly thereafter, which is consistent with the
threshold analysis in Section~\ref{sec:stage2-class}.

\begin{figure}[H]
\includegraphics[width=0.72\linewidth]{precision_recall_curve.pdf}
\caption{Precision--recall curves for both stages on their independent benchmarks. Stage~1
(BRENDA): average precision $=0.844$ against a positive-class baseline of 0.63. Stage~2
(S669): average precision $=0.453$ against a baseline of 0.25. Dotted lines mark the
respective baselines, i.e.\ the precision of a random classifier at that class prevalence.}
\label{fig:pr}
\end{figure}

\begin{figure}[H]
\includegraphics[width=\linewidth]{precision_vs_threshold.pdf}
\caption{Precision--recall trade-off as a function of decision threshold.
(\textbf{a})~Stage~1 (BRENDA): precision rises from 0.65 to 0.98 as the temperature cutoff
tightens from 46 to 66\,$^{\circ}$C, while recall falls from 0.98 to 0.32.
(\textbf{b})~Stage~2 (S669): the ensemble reaches precision 0.81 at the strict
0.00\kcal{} threshold, at 10\% recall. Precision at a strict threshold is not an
experimental success rate.}
\label{fig:precision}
\end{figure}

\subsection{Cross-Stage Comparison}

Stage~1 achieved a higher independent ROC AUC (0.732) than Stage~2 (0.669;
Figure~\ref{fig:roc}). However, the two stages were trained and evaluated on different
datasets, prediction targets (temperature stability vs.\ folding \ddg), class labels, and
sample sizes (1{,}563 vs.\ 669 records). The comparison is therefore descriptive only: it
is not a controlled head-to-head evaluation, and no statistical significance test was
performed between the two stages. The bootstrap intervals do not overlap
(0.708--0.755 vs.\ 0.62--0.72), but because the two quantities are computed on different
data with different targets, that separation should not be read as a significance test.

\begin{figure}[H]
\includegraphics[width=\linewidth]{roc_curve.pdf}
\caption{Independent-benchmark ROC curves for the two-stage framework. (\textbf{a})~Stage~1
BRENDA temperature screening, ROC AUC 0.732 (95\% CI 0.708--0.755) on 1{,}563 proteins.
(\textbf{b})~Stage~2 S669 mutation prioritization, ROC AUC 0.669 (95\% CI 0.62--0.72) on
669 mutations. Dashed line: random classifier (AUC $=0.5$). Positive class: BRENDA $=$
thermostable; S669 $=$ stabilizing.}
\label{fig:roc}
\end{figure}

\subsection{Illustrative Application to a PET Hydrolase}
\label{sec:ispetase}

To illustrate how Stage~1 is applied to a target of interest, we ran the whole-protein model
on the \textit{Ideonella sakaiensis} PET hydrolase (IsPETase; UniProt A0A0K8P6T7,
\texttt{PETH\_PISS1}, 290 residues) \citep{yoshida2016}. This accession does not appear in
the Stage~1 training table, the S669 benchmark, or the curated BRENDA benchmark, so the
prediction is genuinely out of sample.

The three Stage~1 base models produced closely agreeing estimates---LightGBM
55.2\,$^{\circ}$C, XGBoost 55.8\,$^{\circ}$C, CatBoost 55.6\,$^{\circ}$C---for an ensemble
prediction of 55.5\,$^{\circ}$C (Table~\ref{tab:ispetase}). Under the labeling scheme of
Section~\ref{sec:stage1-labels} this places IsPETase in the ambiguous 45--60\,$^{\circ}$C
band rather than in the thermostable class, i.e.\ the model does not nominate the wild-type
enzyme as a thermostable scaffold.

Two cautions apply. First, agreement among base models reflects only their mutual
consistency, not accuracy; Section~\ref{sec:stage1-reg} showed that Stage~1 predictions are
compressed toward this same 45--60\,$^{\circ}$C band, so a prediction inside that band
carries little discriminative information. Second, a single protein cannot validate a
model. This example demonstrates the intended usage pattern; it is not evidence of accuracy
on PET hydrolases, for which no leakage-audited benchmark of comparable size currently
exists.
\todo{a cited experimental \Tm{} for wild-type IsPETase, with reference, so the predicted
55.5\,$^{\circ}$C can be compared against measurement rather than left unanchored}

\begin{table}[H]
\caption{Stage 1 whole-protein prediction for IsPETase (UniProt A0A0K8P6T7, 290 residues).
The accession is absent from the Stage~1 training table, S669, and the curated BRENDA
benchmark.}
\label{tab:ispetase}
\begin{tabular}{lc}
\toprule
\textbf{Model} & \textbf{Predicted \Tm{} ($^{\circ}$C)}\\
\midrule
LightGBM  & 55.2\\
XGBoost   & 55.8\\
CatBoost  & 55.6\\
\midrule
\textbf{Ensemble} & \textbf{55.5}\\
\bottomrule
\end{tabular}
\end{table}

%=================================================================
\section{Discussion}
%=================================================================

\subsection{Interpreting Moderate Performance}

The two-stage framework provides a moderate, leakage-controlled route from whole-protein
screening to mutation-level optimization. Stage~1's independent ROC AUC of 0.732 indicates
useful but imperfect discrimination of thermostable proteins, and its precision--recall
trade-off means the operating point should be chosen to match the intended use (broad
recall for discovery vs.\ high precision for shortlisting).

Stage~2's independent performance (Pearson 0.390; RMSE 1.509\kcal; ROC AUC 0.669) is best
described as moderate, with pronounced high-precision/low-recall behaviour at strict \ddg{}
thresholds. This makes the model a candidate-prioritization filter rather than a
comprehensive predictor of every stabilizing mutation. An RMSE of roughly 1.5\kcal{} is
comparable in magnitude to the effect size of many individually stabilizing mutations,
which is a direct statement of the method's resolution limit: it can enrich a shortlist but
cannot reliably resolve small differences between candidates.

\subsection{Leakage as the Dominant Methodological Risk}

The central methodological contribution is explicit leakage auditing. Removing
approximately 132{,}000 records---the great majority (131{,}479) detectable only through
sequence homology---shows that homology-level contamination, not duplicate records, is the
dominant risk to honest benchmarking in this setting. Exact-key filtering alone would have
caught fewer than 700 of these records.

This has a practical implication beyond our own results: stability benchmarks that report
only exact-match de-duplication cannot be assumed leakage-free, and the gap between the two
filtering levels is large enough that it could plausibly account for differences between
published methods. A residual homology risk remains in our own pipeline---the 30\% identity
and 50\% coverage thresholds are conventional but arbitrary, and remote homology below that
cutoff is not removed---and is acknowledged rather than dismissed.

The internal-versus-independent comparison in Section~\ref{sec:internal-external} makes the
same point from another direction. A development split that respects protein grouping but
draws from the same distribution reported Pearson $\approx$0.80 for Stage~1, against 0.54
on the independent benchmark.

\subsection{Precision Is Not an Experimental Success Rate}

We caution that benchmark precision may not transfer unchanged to PET-degrading enzymes,
which are under-represented in these general stability sets. In particular, a benchmark
precision value such as 0.93 should not be equated with a 93\% experimental success rate; it
is measured at a chosen threshold on curated data, not in the laboratory, and it is
accompanied in our own results by a recall of 0.39.

Reporting precision without its matching recall is the most common way this class of model
is oversold, which is why every precision figure in this paper is stated together with
recall, F1, and accuracy, and why the complete threshold sweeps are provided in
Appendix~\ref{app:sweeps}.

\subsection{Why the Ensemble Does Not Dominate}

The Stage~2 result in which Extra Trees outperforms the unweighted ensemble
(Pearson 0.436 vs.\ 0.390) is worth stating plainly rather than omitting. Unweighted
averaging assumes roughly comparable base-model quality; here the spread was wide
(0.298--0.436), so averaging diluted the strongest member. A validation-weighted ensemble
would likely recover the difference, but weights would have to be fitted on a held-out
split that is not the benchmark, and we did not do so in this study. Selecting Extra Trees
on the basis of its S669 performance would convert an independent benchmark into a model
selection set, which is precisely the practice this paper argues against.

\subsection{Comparison with Prior Work}

Comparisons with prior \Tm{} or \ddg{} predictors are only meaningful when the benchmark,
sign convention, data split, and metric are matched; on S669, a broad set of stability
predictors has been benchmarked under matched conditions \citep{pancotti2022s669}. We
therefore restrict quantitative comparison to genuinely comparable settings and avoid
implying superiority across incomparable evaluations. In particular, we make no claim of
state-of-the-art performance: our Stage~2 correlation is moderate, and the contribution we
claim is the leakage-audited evaluation protocol and the two-stage separation of targets,
not a new accuracy record.

\subsection{Limitations}

Limitations include the moderate correlation on S669, the low recall at high-precision
operating points, the sequence-based (rather than structure-based) scope, residual homology
risk below the 30\% identity threshold, and the fact that both benchmarks are general
stability sets rather than PET-hydrolase-specific assays.

The PET-hydrolase illustration in Section~\ref{sec:ispetase} is a single out-of-sample
prediction and carries no validation weight; no leakage-audited PET-hydrolase benchmark of
comparable size to BRENDA or S669 currently exists, which is the main obstacle to
evaluating the framework on its intended application domain.

The feature representation mean-pools ESM-2 embeddings over the whole sequence, which
attenuates mutation-site signal for Stage~2. Mutations are scored independently, so
epistatic interactions between simultaneous substitutions are not modelled, and multi-mutant
design is outside the framework's current scope. The framework predicts thermodynamic
stability rather than catalytic activity---a stabilized variant may have reduced turnover.
Stage~1 labels derive from heterogeneous BRENDA evidence, only 325 of 2{,}034 curated
proteins carry a directly measured \Tm, and predictions are compressed toward the
45--60\,$^{\circ}$C band. Finally, the software environment was reconstructed from the
training container rather than pinned at run time (Appendix~\ref{app:software}).

\subsection{Future Directions}
\label{sec:future}

Future work should include PET-enzyme-specific evaluation on a leakage-audited benchmark;
per-residue or attention-pooled embeddings to preserve mutation-site detail; a
validation-weighted rather than unweighted ensemble; incorporation of a pH-optimum analysis
once a confirmed training run is available; modelling of epistatic interactions to support
multi-mutant design; and, ultimately, experimental validation before any deployment claim.

%=================================================================
\section{Conclusions}
%=================================================================

We present a two-stage, leakage-audited machine learning framework that separates
whole-protein temperature screening from mutation-level folding-stability optimization,
trained on a harmonized 1{,}001{,}888-record table assembled from seven public sources. On
independent, leakage-audited benchmarks the framework reached ROC AUC 0.732 (95\% CI
0.708--0.755) with average precision 0.844 on BRENDA, and ROC AUC 0.669 (95\% CI 0.62--0.72)
with average precision 0.453 and Pearson 0.390 (95\% CI 0.32--0.46) on S669.

Explicit homology-level auditing removed approximately 132{,}000 records that exact-match
filtering alone would have retained---roughly 197 times as many as the four exact keys
combined---indicating that homology leakage is the dominant contamination risk in this
setting. The gap between internal validation (Pearson $\approx$0.80) and independent
testing (0.54) for Stage~1 quantifies the optimism that an internal split alone would have
reported.

These results describe a moderate, leakage-controlled screen-then-optimize framework
suitable for prioritizing candidates for experimental testing. They do not constitute a
validated industrial solution, and experimental validation remains necessary before any
deployment claim.

%=================================================================
\appendix
\section{Software Environment}
\label{app:software}

Table~\ref{tab:software} lists the pinned environment. These versions were reproduced on
2026-07-25 from the training container using the same install line as the training jobs.
The original training jobs installed these packages unpinned, so patch versions may have
drifted; the metrics should be confirmed against a pinned re-run before these are cited as
definitive.

\begin{table}[H]
\caption{Pinned software environment, reproduced from the training container
\texttt{pytorch/pytorch:2.1.0-cuda12.1-cudnn8-runtime}.}
\label{tab:software}
\begin{tabular}{ll}
\toprule
\textbf{Component} & \textbf{Version}\\
\midrule
Python              & 3.10.13\\
PyTorch             & 2.1.0 (CUDA 12.1)\\
fair-esm            & 2.0.0\\
scikit-learn        & 1.7.2\\
XGBoost             & 3.2.0\\
LightGBM            & 4.7.0\\
CatBoost            & 1.2.10\\
NumPy               & 1.26.0\\
pandas              & 2.3.3\\
joblib              & 1.5.3\\
SciPy               & 1.15.3\\
MMseqs2             & release 18-8cc5c\\
\bottomrule
\end{tabular}
\end{table}

\section{Model Hyperparameters}
\label{app:hyper}

\begin{table}[H]
\caption{Base-model hyperparameters for both stages. All models used random seed 42; the
ensemble prediction is the unweighted mean of the listed base models.}
\label{tab:hyper}
\begin{tabular}{lll}
\toprule
\textbf{Stage} & \textbf{Model} & \textbf{Settings}\\
\midrule
\multirow{5}{*}{Stage 1 (\Tm)}
 & Extra Trees   & 400 trees, max depth 25\\
 & Random Forest & 400 trees, max depth 25\\
 & XGBoost       & 600 rounds, depth 6, learning rate 0.03\\
 & LightGBM      & 800 rounds, learning rate 0.03\\
 & CatBoost      & 800 iterations, depth 8, learning rate 0.03\\
\midrule
\multirow{6}{*}{Stage 2 (\ddg)}
 & Random Forest & 300 trees, max depth 25\\
 & Extra Trees   & 300 trees, max depth 25\\
 & Multilayer perceptron & up to 100 iterations\\
 & XGBoost       & 600 rounds, depth 9, learning rate 0.02, subsample 0.8, GPU hist\\
 & LightGBM      & default rounds, GPU training\\
 & CatBoost      & GPU training\\
\bottomrule
\end{tabular}
\end{table}

\section{Extended Threshold Sweeps}
\label{app:sweeps}

The complete sweeps comprise 1{,}122 distinct thresholds for Stage~1 and 669 for Stage~2.
Tables~\ref{tab:sweep1} and \ref{tab:sweep2} report representative rows spanning the full
range; counts are exact.

\begin{table}[H]
\caption{Stage 1 (BRENDA, $n=1{,}563$; 979 positive, 584 negative): extended
decision-threshold sweep. TP/FP/FN/TN are exact counts at each threshold.}
\label{tab:sweep1}
\begin{tabular}{rcccccccc}
\toprule
\textbf{Cutoff ($^{\circ}$C)} & \textbf{Prec.} & \textbf{Rec.} & \textbf{F1} & \textbf{Acc.} & \textbf{TP} & \textbf{FP} & \textbf{FN} & \textbf{TN}\\
\midrule
41.9 & 0.628 & 1.000 & 0.771 & 0.628 & 979 & 581 & 0   & 3\\
46.0 & 0.650 & 0.977 & 0.781 & 0.656 & 956 & 514 & 23  & 70\\
48.0 & 0.678 & 0.898 & 0.772 & 0.669 & 879 & 418 & 100 & 166\\
50.0 & 0.713 & 0.767 & 0.739 & 0.661 & 751 & 302 & 228 & 282\\
52.0 & 0.766 & 0.629 & 0.691 & 0.647 & 616 & 188 & 363 & 396\\
55.0 & 0.838 & 0.522 & 0.643 & 0.637 & 511 & 99  & 468 & 485\\
58.0 & 0.904 & 0.433 & 0.586 & 0.616 & 424 & 45  & 555 & 539\\
59.9 & 0.926 & 0.394 & 0.553 & 0.601 & 386 & 31  & 593 & 553\\
62.0 & 0.951 & 0.376 & 0.539 & 0.597 & 368 & 19  & 611 & 565\\
65.9 & 0.978 & 0.324 & 0.487 & 0.572 & 317 & 7   & 662 & 577\\
70.1 & 0.989 & 0.287 & 0.445 & 0.552 & 281 & 3   & 698 & 581\\
75.0 & 0.982 & 0.169 & 0.288 & 0.477 & 165 & 3   & 814 & 581\\
80.1 & 0.973 & 0.075 & 0.139 & 0.419 & 73  & 2   & 906 & 582\\
\bottomrule
\end{tabular}
\end{table}

\begin{table}[H]
\caption{Stage 2 (S669, $n=669$; 168 stabilizing, 501 non-stabilizing): extended
decision-threshold sweep. A mutation is predicted stabilizing when its predicted \ddg{}
falls below the cutoff. TP/FP/FN/TN are exact counts.}
\label{tab:sweep2}
\begin{tabular}{rcccccccc}
\toprule
\textbf{Cutoff (\kcal)} & \textbf{Prec.} & \textbf{Rec.} & \textbf{F1} & \textbf{Acc.} & \textbf{TP} & \textbf{FP} & \textbf{FN} & \textbf{TN}\\
\midrule
$-0.52$ & 1.000 & 0.012 & 0.024 & 0.752 & 2   & 0   & 166 & 501\\
$-0.25$ & 0.889 & 0.048 & 0.090 & 0.759 & 8   & 1   & 160 & 500\\
$+0.00$ & 0.818 & 0.107 & 0.189 & 0.770 & 18  & 4   & 150 & 497\\
$+0.25$ & 0.500 & 0.262 & 0.344 & 0.749 & 44  & 44  & 124 & 457\\
$+0.50$ & 0.375 & 0.464 & 0.415 & 0.671 & 78  & 130 & 90  & 371\\
$+0.75$ & 0.331 & 0.655 & 0.440 & 0.581 & 110 & 222 & 58  & 279\\
$+1.00$ & 0.311 & 0.798 & 0.447 & 0.505 & 134 & 297 & 34  & 204\\
$+1.16$ & 0.311 & 0.887 & 0.461 & 0.478 & 149 & 330 & 19  & 171\\
$+1.50$ & 0.278 & 0.917 & 0.427 & 0.381 & 154 & 400 & 14  & 101\\
$+2.00$ & 0.259 & 0.994 & 0.411 & 0.286 & 167 & 477 & 1   & 24\\
\bottomrule
\end{tabular}
\end{table}

%=================================================================
\authorcontributions{(Provisional --- authors to confirm.) A.I.\ contributed to
conceptualization, software, data curation, formal analysis, investigation, methodology,
visualization, and writing of the original draft. J.P.\ contributed to conceptualization,
supervision, validation, and writing (review and editing). Ab.I.\ contributed to
investigation, writing of the original draft, and writing (review and editing). All authors
reviewed and approved the final manuscript.}
% CONFIRM/ADJUST the CRediT roles above and the author order before submission.

\funding{No specific grant funding was received for this work; computational resources were
provided by the National Research Platform (Nautilus).}
% CONFIRM funding status with all authors (change if a grant supported the work).

\dataavailability{The processed training table, leakage-audited benchmark datasets, saved
predictions, split identifiers, leakage-audit outputs, and trained model artifacts
generated in this study are available through Zenodo
(\url{https://doi.org/10.5281/zenodo.21257369}); source datasets are cited above. Source
code for data assembly, leakage auditing, model training, evaluation, threshold analysis,
bootstrap confidence-interval estimation, and figure generation is available at
\url{https://github.com/AyushIyer31/PET-Lab} and archived on Zenodo at
\url{https://doi.org/10.5281/zenodo.21519961} (release v1.0.0).}
% ACTION before submission: make the Zenodo record open (or provide reviewer access) and
% apply a data license such as CC BY 4.0 instead of the current GPL-3.0; confirm the
% deposit contains all listed artifacts with a README/data dictionary. CONFIRM the archived
% commit (497c4f5) is the exact code that produced the results.

\acknowledgments{The authors thank the National Research Platform (Nautilus) for providing
computational resources. Nautilus is supported by the Pacific Research Platform (NSF
award~1541349), CHASE-CI (NSF award~1730158), and Towards a National Research Platform (NSF
award~1826967), with additional funding from the University of California Office of the
President.}
% CONFIRM: verify these NRP/NSF award numbers are current for your usage, and add any
% individual (non-author) acknowledgements with their consent.

\conflictsofinterest{The authors declare no conflicts of interest.}
% CONFIRM with all authors before submission.

% Large-language-model use: a large language model (Anthropic Claude) was used to assist
% with drafting, formatting, and reference preparation; all content was reviewed and
% verified by the authors, who take full responsibility for it.
% CONFIRM wording against the journal's current LLM-use policy and place it in the section
% the journal requires.

%=================================================================
\begin{thebibliography}{999}

\bibitem[OECD(2022)]{oecd2022plastics}
OECD. \textit{Global Plastics Outlook: Economic Drivers, Environmental Impacts and Policy
Options}; OECD Publishing: Paris, \textbf{2022}. doi:10.1787/de747aef-en.

\bibitem[Yoshida et al.(2016)]{yoshida2016}
Yoshida, S.; Hiraga, K.; Takehana, T.; Taniguchi, I.; Yamaji, H.; Maeda, Y.; Toyohara, K.;
Miyamoto, K.; Kimura, Y.; Oda, K. A bacterium that degrades and assimilates
poly(ethylene terephthalate). \textit{Science} \textbf{2016}, \textit{351}, 1196--1199.
doi:10.1126/science.aad6359.

\bibitem[Tournier et al.(2020)]{tournier2020}
Tournier, V.; Topham, C.M.; Gilles, A.; David, B.; Folgoas, C.; Moya-Leclair, E.; Kamionka,
E.; Desrousseaux, M.-L.; Texier, H.; Gavalda, S.; et al. An engineered PET depolymerase to
break down and recycle plastic bottles. \textit{Nature} \textbf{2020}, \textit{580},
216--219. doi:10.1038/s41586-020-2149-4.

\bibitem[Lu et al.(2022)]{lu2022}
Lu, H.; Diaz, D.J.; Czarnecki, N.J.; Zhu, C.; Kim, W.; Shroff, R.; Acosta, D.J.; Alexander,
B.R.; Cole, H.O.; Zhang, Y.; et al. Machine learning-aided engineering of hydrolases for
PET depolymerization. \textit{Nature} \textbf{2022}, \textit{604}, 662--667.
doi:10.1038/s41586-022-04599-z.

\bibitem[Tsuboyama et al.(2023)]{tsuboyama2023}
Tsuboyama, K.; Dauparas, J.; Chen, J.; Laine, E.; Mohseni Behbahani, Y.; Weinstein, J.J.;
Mangan, N.M.; Ovchinnikov, S.; Rocklin, G.J. Mega-scale experimental analysis of protein
folding stability in biology and design. \textit{Nature} \textbf{2023}, \textit{620},
434--444. doi:10.1038/s41586-023-06328-6.

\bibitem[Jarzab et al.(2020)]{jarzab2020meltome}
Jarzab, A.; Kurzawa, N.; Hopf, T.; Moerch, M.; Zecha, J.; Leijten, N.; Bian, Y.; Musiol, E.;
Maschberger, M.; Stoehr, G.; et al. Meltome atlas---thermal proteome stability across the
tree of life. \textit{Nature Methods} \textbf{2020}, \textit{17}, 495--503.
doi:10.1038/s41592-020-0801-4.

\bibitem[Xavier et al.(2021)]{xavier2021thermomutdb}
Xavier, J.S.; Nguyen, T.-B.; Karmarkar, M.; Portelli, S.; Rezende, P.M.; Velloso, J.P.L.;
Ascher, D.B.; Pires, D.E.V. ThermoMutDB: a thermodynamic database for missense mutations.
\textit{Nucleic Acids Research} \textbf{2021}, \textit{49}, D475--D479.
doi:10.1093/nar/gkaa925.

\bibitem[Stourac et al.(2021)]{stourac2021fireprotdb}
Stourac, J.; Dubrava, J.; Musil, M.; Horackova, J.; Damborsky, J.; Mazurenko, S.; Bednar, D.
FireProtDB: database of manually curated protein stability data. \textit{Nucleic Acids
Research} \textbf{2021}, \textit{49}, D319--D324. doi:10.1093/nar/gkaa981.

\bibitem[Beltran et al.(2025)]{beltran2025domainome}
Beltran, A.; Jiang, X.; Shen, Y.; Lehner, B. Site-saturation mutagenesis of 500 human
protein domains. \textit{Nature} \textbf{2025}. doi:10.1038/s41586-024-08370-4.
% VERIFY volume/pages at proof stage.

\bibitem[Pancotti et al.(2022)]{pancotti2022s669}
Pancotti, C.; Benevenuta, S.; Birolo, G.; Alberini, V.; Repetto, V.; Sanavia, T.; Capriotti,
E.; Fariselli, P. Predicting protein stability changes upon single-point mutation: a
thorough comparison of the available tools on a new dataset. \textit{Briefings in
Bioinformatics} \textbf{2022}, \textit{23}, bbab555. doi:10.1093/bib/bbab555.

\bibitem[Dehouck et al.(2009)]{dehouck2009s2648}
Dehouck, Y.; Grosfils, A.; Folch, B.; Gilis, D.; Bogaerts, P.; Rooman, M. Fast and accurate
predictions of protein stability changes upon mutations using statistical potentials and
neural networks: PoPMuSiC-2.0. \textit{Bioinformatics} \textbf{2009}, \textit{25},
2537--2543. doi:10.1093/bioinformatics/btp445.
% Source of the S2648 dataset; VERIFY that this is the ProDDG/S2648 provenance used.

\bibitem[Chang et al.(2021)]{chang2021brenda}
Chang, A.; Jeske, L.; Ulbrich, S.; Hofmann, J.; Koblitz, J.; Schomburg, I.; Neumann-Schaal,
M.; Jahn, D.; Schomburg, D. BRENDA, the ELIXIR core data resource in 2021: new developments
and updates. \textit{Nucleic Acids Research} \textbf{2021}, \textit{49}, D498--D508.
doi:10.1093/nar/gkaa1025.

\bibitem[Steinegger and S\"{o}ding(2017)]{steinegger2017mmseqs2}
Steinegger, M.; S\"{o}ding, J. MMseqs2 enables sensitive protein sequence searching for the
analysis of massive data sets. \textit{Nature Biotechnology} \textbf{2017}, \textit{35},
1026--1028. doi:10.1038/nbt.3988.

\bibitem[Lin et al.(2023)]{lin2023esm2}
Lin, Z.; Akin, H.; Rao, R.; Hie, B.; Zhu, Z.; Lu, W.; Smetanin, N.; Verkuil, R.; Kabeli, O.;
Shmueli, Y.; et al. Evolutionary-scale prediction of atomic-level protein structure with a
language model. \textit{Science} \textbf{2023}, \textit{379}, 1123--1130.
doi:10.1126/science.ade2574.

\bibitem[Pedregosa et al.(2011)]{pedregosa2011sklearn}
Pedregosa, F.; Varoquaux, G.; Gramfort, A.; Michel, V.; Thirion, B.; Grisel, O.; Blondel, M.;
Prettenhofer, P.; Weiss, R.; Dubourg, V.; et al. Scikit-learn: Machine learning in Python.
\textit{Journal of Machine Learning Research} \textbf{2011}, \textit{12}, 2825--2830.

\bibitem[Chen and Guestrin(2016)]{chen2016xgboost}
Chen, T.; Guestrin, C. XGBoost: A scalable tree boosting system. In \textit{Proceedings of
the 22nd ACM SIGKDD International Conference on Knowledge Discovery and Data Mining},
\textbf{2016}; pp. 785--794. doi:10.1145/2939672.2939785.

\bibitem[Ke et al.(2017)]{ke2017lightgbm}
Ke, G.; Meng, Q.; Finley, T.; Wang, T.; Chen, W.; Ma, W.; Ye, Q.; Liu, T.-Y. LightGBM: A
highly efficient gradient boosting decision tree. In \textit{Advances in Neural Information
Processing Systems}, \textbf{2017}; Volume 30.

\bibitem[Prokhorenkova et al.(2018)]{prokhorenkova2018catboost}
Prokhorenkova, L.; Gusev, G.; Vorobev, A.; Dorogush, A.V.; Gulin, A. CatBoost: unbiased
boosting with categorical features. In \textit{Advances in Neural Information Processing
Systems}, \textbf{2018}; Volume 31.

\end{thebibliography}

\end{document}
